\chapter{Implementation}
This section provides an overview of what envoriment was used to implement the
models, what libraries were used and why, how the classes were designed.
\section{Coding envoriment}
The implementation of the HMM models is in Python 3.11 and is split into two
parts: a model library, which contains the HMMs themselves, and a set of
scripts that clean the data, fit the models and evaluate them. The model
library only depends on a small numerical stack:
\begin{itemize}
    \item \textbf{JAX} (0.9.0) \parencite{jax2018}. All model computations,
    i.e.\ the emission densities, the transition matrices and the forward
    algorithm, are written in JAX. JAX provides automatic differentiation, so
    the gradient of the log-likelihood with respect to every model parameter is
    obtained without deriving it by hand. This is what allows the models to be
    fitted by direct numerical optimisation of the likelihood instead of the
    classical Baum-Welch (EM) algorithm, and it makes it cheap to add new
    emission or transition types, since no new update equations have to be
    derived. JAX also just-in-time compiles the code, and the recursion over
    time in the forward algorithm is written with \texttt{jax.lax.scan}, which
    keeps the fitting fast. Since the likelihood is numerically sensitive, all
    computations are done in 64-bit floating point.
    \item \textbf{Equinox} (0.13.7) \parencite{kidger2021equinox}. Equinox lets
    the model classes be written as ordinary Python classes that are at the same
    time JAX pytrees. An HMM is therefore a single object whose parameters JAX
    can differentiate and optimise directly, and Equinox's filtering utilities
    are used to freeze selected parameters during a fit.
    \item \textbf{Optax} (0.2.8) \parencite{deepmind2020jax}. Optax provides the
    optimisers that minimise the negative log-likelihood. The default is the
    quasi-Newton method L-BFGS, while first-order gradient methods such as Adam
    are available as an alternative.
    \item \textbf{NumPy} (2.2) \parencite{harris2020numpy} for array handling
    outside of the JAX code.
\end{itemize}
The scripts around the library additionally use the following libraries:
\begin{itemize}
    \item \textbf{pandas} (3.0) for loading, cleaning and aggregating the
    $\text{CO}_2$ time series into 30 minute intervals.
    \item \textbf{Matplotlib} (3.10) and \textbf{seaborn} (0.13) for all
    figures.
    \item \textbf{statsmodels} (0.14) for the autocorrelation functions and
    QQ-plots used to evaluate the pseudo-residuals of the fitted models.
    \item \textbf{holidays} (0.105) to mark Danish public holidays when
    constructing the time-of-day and weekday covariates.
    \item \textbf{Requests} (2.34) and \textbf{pyproj} (3.7) to download outdoor
    humidity data from the open data API of the Danish Meteorological Institute
    (DMI), where pyproj converts the coordinates of the building into the
    coordinate system used by DMI.
    \item \textbf{python-dotenv} (1.2) to read the location of the data from a
    local configuration file, so that the data itself is kept out of the code
    repository.
\end{itemize}
The package manager used is uv \parencite{uv}. The exact versions of every
library are recorded in a lock file, so the environment can be reproduced with
a single \texttt{uv sync} command, and uv also installs the pinned Python
version. This ensures that the results in this report can be reproduced with
the same versions of the libraries that produced them.

\section{Class Design}
The class design of the model library is shown in \cref{fig:class_design}.
An HMM is described by a \texttt{HMMParams} object, which is composed of a
transition model and an emission model. Both are given as abstract base classes,
\texttt{BaseTransition} and \texttt{BaseEmission}, with the concrete models
from the theory chapter implemented as subclasses. Because the HMM is
assembled from these two parts, a new model is obtained by combining an existing
transition with an existing emission, e.g.\ a covariate-driven transition with
an autoregressive emission, without writing any new code for the combination.

The remaining classes act on such a parameter object. The inference package
contains the forward algorithm, which is used both by the negative
log-likelihood that the solvers minimise, and by the diagnostics and the
predictions on a fitted model. The solvers, the inference algorithm and the
prediction are each given by an abstract base class, so an alternative
optimiser or inference algorithm can be substituted without changing the model
classes.

\begin{figure}[H]
    \centering
    % UML class diagram of the model library (fig:class_design).
% Needs the TikZ libraries positioning, fit, backgrounds, shapes.multipart and
% arrows.meta, and the colours umlfill, umlborder, pkgborder and pkgtext (all
% in Setup/Settings.tex). Font sizes are fixed so the layout does not depend on
% the document's base size; it is drawn to fit a 15 cm text width, so do not
% scale it.
\begin{tikzpicture}[
    font=\fontsize{8}{9.6}\selectfont,
    % Class with a name compartment and a member compartment.
    cls/.style={draw=umlborder, fill=umlfill, line width=0.6pt,
        rectangle split, rectangle split parts=2,
        rectangle split part align={center, left},
        inner xsep=5pt, inner ysep=4pt},
    % Class with only a name compartment.
    leaf/.style={draw=umlborder, fill=umlfill, line width=0.6pt,
        inner xsep=5pt, inner ysep=4pt},
    % Concrete class: fixed text height so boxes in a row line up.
    concrete/.style={leaf, text height=1.6ex, text depth=0.4ex},
    pkg/.style={draw=pkgborder, line width=0.6pt, inner xsep=6pt, inner ysep=6pt},
    pkgname/.style={anchor=north west, font=\fontsize{9.5}{11}\selectfont\bfseries, text=pkgtext,
        inner sep=5pt},
    line/.style={umlborder, line width=0.7pt},
    inherit/.style={line, -{Triangle[open, length=8pt, width=9pt]}},
    realize/.style={inherit, dashed},
    uses/.style={line, -{Stealth[length=7pt]}},
    depends/.style={uses, dashed},
    composed/.style={line, {Diamond[length=10pt, width=7pt]}-},
    aggregates/.style={line, {Diamond[open, length=10pt, width=7pt]}-},
    lbl/.style={ text=black, fill=white, inner sep=2pt},
]
\def\namefont{\fontsize{9}{11}\selectfont}
% \name{<stereotype>}{<name>} and \members{<line>\\<line>...}.
\def\name#1#2{\begin{tabular}{@{}c@{}}\fontsize{6.5}{8}\selectfont\guillemotleft#1\guillemotright\\\namefont#2\end{tabular}}
\def\members#1{\begin{tabular}{@{}l@{}}#1\end{tabular}}

% ================================================================ fitting
\node[cls, anchor=north west] (loss) at (0.3, 0) {
    \name{interface}{LossFunction}
    \nodepart{second}\members{(params, data) $\to$ float}};
\node[leaf, below=0.6cm of loss] (nll) {\name{function}{negative\_log\_likelihood}};
% Common left and right edge for the packages.
\coordinate (L) at (nll.west);
\coordinate (R) at (14.1, 0);
\node[cls, anchor=north east] (solver) at (14.1, 0) {
    \name{abstract}{\textit{BaseSolver}}
    \nodepart{second}\members{+ fit(params, loss\_fn, data) : HMMParams}};
\node[concrete, anchor=north east] (grad) at ([yshift=-0.6cm] solver.south east) {GradientSolver};
\node[concrete, anchor=north west] (lbfgs) at ([yshift=-0.6cm] solver.south west) {LBFGSSolver};

\draw[inherit] (lbfgs) -- (lbfgs |- solver.south);
\draw[inherit] (grad) -- (grad |- solver.south);
\draw[uses] (solver.west) -- node[lbl, above=2pt] {minimises} (solver.west -| loss.east);
\draw[realize] (nll) -- (loss);

% ================================================================ inference
\node[cls, anchor=north west] (forward) at ([yshift=-1.15cm] nll.south -| 0.3, 0) {
    \namefont ForwardAlgorithm
    \nodepart{second}\members{+ run(params, data) : ForwardOutput}};
\node[cls, anchor=north east] (inference) at (14.1, 0 |- forward.north) {
    \name{abstract}{\textit{BaseInference}}
    \nodepart{second}\members{+ run(params, data) : BaseOutput}};

\draw[inherit] (forward.east |- inference.text split) -- (inference.west |- inference.text split);
\draw[depends] (nll) -- node[lbl, pos=0.37, right=2pt] {calls} (nll |- forward.north);

% ================================================================ data
\node[cls, anchor=north west] (data) at ([yshift=-1.15cm] inference.south -| 0.3, 0) {
    \name{pydantic dataclass}{Data}
    \nodepart{second}\members{+ ys : Array\\
        + ts : Array | None\\
        + xs : Array | None\\
        + initial\_distribution : Array}};

% ================================================================ analysis
\node[cls, anchor=north east] (diagnostics) at (14.1, 0 |- data.north) {
    \namefont Diagnostics
    \nodepart{second}\members{+ params : HMMParams\\
        + data : Data\\
        + inference : BaseInference\\[2pt]\hline\rule{0pt}{10pt}%
        + log\_likelihood() : float\\
        + state\_results() : StateResults\\
        + pseudo\_residuals() : Array}};
\node[cls, anchor=north east] (prediction) at ([xshift=-0.9cm] diagnostics.north west) {
    \name{abstract}{\textit{BasePrediction}}
    \nodepart{second}\members{+ predict(params, data,\\
        \hspace{2em}t\_pred, x\_pred) : Array}};
\node[cls, below=0.6cm of prediction] (meanpred) {
    \namefont EmissionMeanPrediction
    \nodepart{second}\members{+ inference : BaseInference}};

\draw[inherit] (meanpred) -- (prediction);
\draw[aggregates] (diagnostics.north -| inference.south) -- (inference.south);
\draw[aggregates] (meanpred.east) -- ++(0.45, 0)
    |- ([yshift=-0.6cm] inference.south west)
    -- ++(0.6, 0) -- ([xshift=0.6cm] inference.south west);

% ================================================================ model
\node[cls] (hmm) at ([yshift=-1.85cm] meanpred.south -| 7.1, 0) {
    \namefont HMMParams
    \nodepart{second}\members{+ transition : BaseTransition\\
        + emission : BaseEmission}};
\node[cls, anchor=north] (transition) at ([yshift=-0.7cm] hmm.south -| 3.5, 0) {
    \name{abstract}{\textit{BaseTransition}}
    \nodepart{second}\members{+ transition\_matrix(t, ys, xs, dt) : Array}};
\node[cls, anchor=north] (emission) at ([yshift=-0.7cm] hmm.south -| 11.6, 0) {
    \name{abstract}{\textit{BaseEmission}}
    \nodepart{second}\members{+ density(t, ys, xs, ts) : Array\\
        + cdf(t, ys, xs, ts) : Array\\
        + mu(t, ys, xs, ts) : Array}};

\draw[composed] (hmm.west |- hmm.text split) -| (transition.north);
\draw[composed] (hmm.east |- hmm.text split) -| (emission.north);

% Concrete classes hang off a vertical trunk to the left of each column.
\coordinate (trunkA) at (0.35, 0);
\coordinate (trunkB) at (4.95, 0);
\coordinate (trunkC) at (9.65, 0);
\node[concrete, anchor=north west] (static) at ([xshift=0.3cm, yshift=-0.9cm] trunkA |- emission.south) {StaticTransition};
\node[concrete, below=0.4cm of static.south west, anchor=north west] (cstatic) {ContinuousStaticTransition};
\node[concrete, anchor=north west] (dynamic) at ([xshift=0.3cm] trunkB |- static.north) {DynamicTransition};
\node[concrete, below=0.4cm of dynamic.south west, anchor=north west] (cdynamic) {ContinuousDynamicTransition};
\node[concrete, anchor=north west] (gauss) at ([xshift=0.3cm] trunkC |- static.north) {GaussEmission};
\node[concrete, below=0.4cm of gauss.south west, anchor=north west] (argauss) {AutoregressiveGaussEmission};

\coordinate (tbus) at ([yshift=0.55cm] static.north -| transition.south);
\draw[inherit] (tbus) -- (transition.south);
\draw[line] (trunkA |- cstatic) |- (tbus) -| (trunkB |- cdynamic);
\foreach \c in {static, cstatic} \draw[line] (trunkA |- \c) -- (\c.west);
\foreach \c in {dynamic, cdynamic} \draw[line] (trunkB |- \c) -- (\c.west);
\coordinate (ebus) at ([yshift=0.55cm] gauss.north -| emission.south);
\draw[inherit] (ebus) -- (emission.south);
\draw[line] (trunkC |- argauss) |- (ebus);
\foreach \c in {gauss, argauss} \draw[line] (trunkC |- \c) -- (\c.west);

% ================================================================ packages
\begin{scope}[on background layer]
    \node[pkg, fit={(loss)(nll)(solver)(lbfgs)(grad)(L)(R)([yshift=16pt] loss.north)}] (pFit) {};
    \node[pkg, fit={(forward)(inference)(L |- forward)(R |- forward)([yshift=16pt] inference.north)}] (pInf) {};
    \node[pkg, fit={(data)(L |- data)([yshift=16pt] data.north)}] (pData) {};
    \node[pkg, fit={(prediction)(meanpred)(diagnostics)([yshift=16pt] diagnostics.north)}] (pAna) {};
    \node[pkg, fit={(hmm)(transition)(emission)(cstatic)(cdynamic)(argauss)(L |- static)(R |- static)([yshift=16pt] hmm.north)}] (pMod) {};
\end{scope}
\node[pkgname] at (pFit.north west) {Fitting};
\node[pkgname] at (pInf.north west) {Inference};
\node[pkgname] at (pData.north west) {Data};
\node[pkgname] at (pAna.north west) {Analysis};
\node[pkgname] at (pMod.north west) {Model};
\end{tikzpicture}

    \caption{Class diagram of the model library, grouped into packages.
    Arrows with an open triangle denote inheritance (dashed: implementation of
    an interface), filled diamonds composition, open diamonds aggregation, and
    dashed arrows a call. Abstract classes are written in italics.}
    \label{fig:class_design}
\end{figure}

